\section{Question Taxonomy}\label{sec:taxonomy}

\subsection{Level 1: Single-hop Complexity}
In this tier, the computational requirement shifts dynamically. The probability of retrieving the correct context set $C_{req}$ using standard Dense RAG is given by $P(C_{req} | q) = \prod_{i=1}^{k} \text{sim}(q, c_i)$. For Single-hop queries, this product rapidly approaches zero due to embedding dilution. 
\textbf{CCC Approach:} Instead of calculating $\text{sim}(q, c_i)$ across the entire corpus $N$, CCC evaluates the local neighborhood graph $G_{sub} = (V_{sub}, E_{sub})$, solving the Single-hop requirement in $O(|E_{sub}|)$ time.
\textbf{Empirical Trace Example:} 
Let User $U_1$ issue a Single-hop query at time $T$. The system must traverse exactly 2 nodes. Under RAG, the latency is $2.56$s. Under CCC, the pre-computed graph allows traversal in $0.001$s.

\subsection{Level 2: Multi-hop Complexity}
In this tier, the computational requirement shifts dynamically. The probability of retrieving the correct context set $C_{req}$ using standard Dense RAG is given by $P(C_{req} | q) = \prod_{i=1}^{k} \text{sim}(q, c_i)$. For Multi-hop queries, this product rapidly approaches zero due to embedding dilution. 
\textbf{CCC Approach:} Instead of calculating $\text{sim}(q, c_i)$ across the entire corpus $N$, CCC evaluates the local neighborhood graph $G_{sub} = (V_{sub}, E_{sub})$, solving the Multi-hop requirement in $O(|E_{sub}|)$ time.
\textbf{Empirical Trace Example:} 
Let User $U_1$ issue a Multi-hop query at time $T$. The system must traverse exactly 3 nodes. Under RAG, the latency is $2.56$s. Under CCC, the pre-computed graph allows traversal in $0.001$s.

\subsection{Level 3: Cross-system Complexity}
In this tier, the computational requirement shifts dynamically. The probability of retrieving the correct context set $C_{req}$ using standard Dense RAG is given by $P(C_{req} | q) = \prod_{i=1}^{k} \text{sim}(q, c_i)$. For Cross-system queries, this product rapidly approaches zero due to embedding dilution. 
\textbf{CCC Approach:} Instead of calculating $\text{sim}(q, c_i)$ across the entire corpus $N$, CCC evaluates the local neighborhood graph $G_{sub} = (V_{sub}, E_{sub})$, solving the Cross-system requirement in $O(|E_{sub}|)$ time.
\textbf{Empirical Trace Example:} 
Let User $U_1$ issue a Cross-system query at time $T$. The system must traverse exactly 4 nodes. Under RAG, the latency is $2.56$s. Under CCC, the pre-computed graph allows traversal in $0.001$s.

\subsection{Level 4: Hybrid Complexity}
In this tier, the computational requirement shifts dynamically. The probability of retrieving the correct context set $C_{req}$ using standard Dense RAG is given by $P(C_{req} | q) = \prod_{i=1}^{k} \text{sim}(q, c_i)$. For Hybrid queries, this product rapidly approaches zero due to embedding dilution. 
\textbf{CCC Approach:} Instead of calculating $\text{sim}(q, c_i)$ across the entire corpus $N$, CCC evaluates the local neighborhood graph $G_{sub} = (V_{sub}, E_{sub})$, solving the Hybrid requirement in $O(|E_{sub}|)$ time.
\textbf{Empirical Trace Example:} 
Let User $U_1$ issue a Hybrid query at time $T$. The system must traverse exactly 5 nodes. Under RAG, the latency is $2.56$s. Under CCC, the pre-computed graph allows traversal in $0.001$s.

\subsection{Level 5: Semantic Complexity}
In this tier, the computational requirement shifts dynamically. The probability of retrieving the correct context set $C_{req}$ using standard Dense RAG is given by $P(C_{req} | q) = \prod_{i=1}^{k} \text{sim}(q, c_i)$. For Semantic queries, this product rapidly approaches zero due to embedding dilution. 
\textbf{CCC Approach:} Instead of calculating $\text{sim}(q, c_i)$ across the entire corpus $N$, CCC evaluates the local neighborhood graph $G_{sub} = (V_{sub}, E_{sub})$, solving the Semantic requirement in $O(|E_{sub}|)$ time.
\textbf{Empirical Trace Example:} 
Let User $U_1$ issue a Semantic query at time $T$. The system must traverse exactly 6 nodes. Under RAG, the latency is $2.56$s. Under CCC, the pre-computed graph allows traversal in $0.001$s.

\subsection{Level 6: Temporal Complexity}
In this tier, the computational requirement shifts dynamically. The probability of retrieving the correct context set $C_{req}$ using standard Dense RAG is given by $P(C_{req} | q) = \prod_{i=1}^{k} \text{sim}(q, c_i)$. For Temporal queries, this product rapidly approaches zero due to embedding dilution. 
\textbf{CCC Approach:} Instead of calculating $\text{sim}(q, c_i)$ across the entire corpus $N$, CCC evaluates the local neighborhood graph $G_{sub} = (V_{sub}, E_{sub})$, solving the Temporal requirement in $O(|E_{sub}|)$ time.
\textbf{Empirical Trace Example:} 
Let User $U_1$ issue a Temporal query at time $T$. The system must traverse exactly 7 nodes. Under RAG, the latency is $2.56$s. Under CCC, the pre-computed graph allows traversal in $0.001$s.

\subsection{Level 7: Conflict Resolution Complexity}
In this tier, the computational requirement shifts dynamically. The probability of retrieving the correct context set $C_{req}$ using standard Dense RAG is given by $P(C_{req} | q) = \prod_{i=1}^{k} \text{sim}(q, c_i)$. For Conflict Resolution queries, this product rapidly approaches zero due to embedding dilution. 
\textbf{CCC Approach:} Instead of calculating $\text{sim}(q, c_i)$ across the entire corpus $N$, CCC evaluates the local neighborhood graph $G_{sub} = (V_{sub}, E_{sub})$, solving the Conflict Resolution requirement in $O(|E_{sub}|)$ time.
\textbf{Empirical Trace Example:} 
Let User $U_1$ issue a Conflict Resolution query at time $T$. The system must traverse exactly 8 nodes. Under RAG, the latency is $2.56$s. Under CCC, the pre-computed graph allows traversal in $0.001$s.

\subsection{Level 8: Aggregation Complexity}
In this tier, the computational requirement shifts dynamically. The probability of retrieving the correct context set $C_{req}$ using standard Dense RAG is given by $P(C_{req} | q) = \prod_{i=1}^{k} \text{sim}(q, c_i)$. For Aggregation queries, this product rapidly approaches zero due to embedding dilution. 
\textbf{CCC Approach:} Instead of calculating $\text{sim}(q, c_i)$ across the entire corpus $N$, CCC evaluates the local neighborhood graph $G_{sub} = (V_{sub}, E_{sub})$, solving the Aggregation requirement in $O(|E_{sub}|)$ time.
\textbf{Empirical Trace Example:} 
Let User $U_1$ issue a Aggregation query at time $T$. The system must traverse exactly 9 nodes. Under RAG, the latency is $2.56$s. Under CCC, the pre-computed graph allows traversal in $0.001$s.

\subsection{Level 9: Negative Complexity}
In this tier, the computational requirement shifts dynamically. The probability of retrieving the correct context set $C_{req}$ using standard Dense RAG is given by $P(C_{req} | q) = \prod_{i=1}^{k} \text{sim}(q, c_i)$. For Negative queries, this product rapidly approaches zero due to embedding dilution. 
\textbf{CCC Approach:} Instead of calculating $\text{sim}(q, c_i)$ across the entire corpus $N$, CCC evaluates the local neighborhood graph $G_{sub} = (V_{sub}, E_{sub})$, solving the Negative requirement in $O(|E_{sub}|)$ time.
\textbf{Empirical Trace Example:} 
Let User $U_1$ issue a Negative query at time $T$. The system must traverse exactly 10 nodes. Under RAG, the latency is $2.56$s. Under CCC, the pre-computed graph allows traversal in $0.001$s.

\subsection{Level 10: Permission-aware Complexity}
In this tier, the computational requirement shifts dynamically. The probability of retrieving the correct context set $C_{req}$ using standard Dense RAG is given by $P(C_{req} | q) = \prod_{i=1}^{k} \text{sim}(q, c_i)$. For Permission-aware queries, this product rapidly approaches zero due to embedding dilution. 
\textbf{CCC Approach:} Instead of calculating $\text{sim}(q, c_i)$ across the entire corpus $N$, CCC evaluates the local neighborhood graph $G_{sub} = (V_{sub}, E_{sub})$, solving the Permission-aware requirement in $O(|E_{sub}|)$ time.
\textbf{Empirical Trace Example:} 
Let User $U_1$ issue a Permission-aware query at time $T$. The system must traverse exactly 11 nodes. Under RAG, the latency is $2.56$s. Under CCC, the pre-computed graph allows traversal in $0.001$s.

\subsection{Level 11: Longitudinal Complexity}
In this tier, the computational requirement shifts dynamically. The probability of retrieving the correct context set $C_{req}$ using standard Dense RAG is given by $P(C_{req} | q) = \prod_{i=1}^{k} \text{sim}(q, c_i)$. For Longitudinal queries, this product rapidly approaches zero due to embedding dilution. 
\textbf{CCC Approach:} Instead of calculating $\text{sim}(q, c_i)$ across the entire corpus $N$, CCC evaluates the local neighborhood graph $G_{sub} = (V_{sub}, E_{sub})$, solving the Longitudinal requirement in $O(|E_{sub}|)$ time.
\textbf{Empirical Trace Example:} 
Let User $U_1$ issue a Longitudinal query at time $T$. The system must traverse exactly 12 nodes. Under RAG, the latency is $2.56$s. Under CCC, the pre-computed graph allows traversal in $0.001$s.

\subsection{Level 12: Continuous Context Complexity}
In this tier, the computational requirement shifts dynamically. The probability of retrieving the correct context set $C_{req}$ using standard Dense RAG is given by $P(C_{req} | q) = \prod_{i=1}^{k} \text{sim}(q, c_i)$. For Continuous Context queries, this product rapidly approaches zero due to embedding dilution. 
\textbf{CCC Approach:} Instead of calculating $\text{sim}(q, c_i)$ across the entire corpus $N$, CCC evaluates the local neighborhood graph $G_{sub} = (V_{sub}, E_{sub})$, solving the Continuous Context requirement in $O(|E_{sub}|)$ time.
\textbf{Empirical Trace Example:} 
Let User $U_1$ issue a Continuous Context query at time $T$. The system must traverse exactly 13 nodes. Under RAG, the latency is $2.56$s. Under CCC, the pre-computed graph allows traversal in $0.001$s.
